\section*{Exercises}
\addcontentsline{toc}{section}{Exercises}

\subsection*{A. Theory}

\begin{exercise}
\exerciseid{16.A1}
Explain why a low-rank signal plus full-rank noise may still have a visible
singular-value elbow.
\end{exercise}

\begin{exercise}
\exerciseid{16.A2}
Explain why a rank-$k$ factorization $LR$ uses far fewer parameters than a full
$m\times n$ matrix when $k\ll \min(m,n)$.
\end{exercise}

\subsection*{B. By Hand}

\begin{exercise}
\exerciseid{16.B1}
Suppose a matrix has singular values $10,7,1,0.8,0.7$. Which rank would you try
first for denoising, and why?
\end{exercise}

\begin{exercise}
\exerciseid{16.B2}
A $100\times 80$ rank-$3$ factorization $LR$ uses
$L\in\R^{100\times 3}$ and $R\in\R^{3\times 80}$. Compare the number of
parameters with the full $100\times 80$ matrix.
\end{exercise}

\subsection*{C. Python / NumPy}

\begin{exercise}
\exerciseid{16.C1}
Generate a rank-$2$ matrix plus Gaussian noise. Compute singular values and
compare rank-$1$, rank-$2$, and rank-$5$ reconstructions.
\end{exercise}

\begin{exercise}
\exerciseid{16.C2}
For your noisy matrix from 16.C1, compute the energy fraction kept by the top
$k$ singular values for $k=1,2,5$.
\end{exercise}

\subsection*{D. Challenge}

\begin{exercise}
\exerciseid{16.D1}
Explain why too small a truncation rank underfits while too large a rank leaves
noise in the reconstruction.
\end{exercise}
